% ═══════════════════════════════════════════════════════════════
%  Chapter 1 — General Context of the Project
% ═══════════════════════════════════════════════════════════════
\chapter{General Context of the Project}
\label{chap:context}
\section*{Chapter Introduction}

This first chapter situates the internship within its professional and technical context. It begins with a presentation of the host organizations---IBM CIC Morocco and Hakkoda---and introduces the client company, ARYZTA, along with the role of Snowflake in its data ecosystem. The chapter then proceeds to articulate the problem addressed by the project, outlining the existing environment architecture, the discrepancies observed, and the risks they pose to ARYZTA's data operations. It concludes with a statement of the project's objectives and a description of the Agile management methodology employed throughout the internship.

% ─────────────────────────────────────────────────────────────
\section{General Overview of the Host Organization}
\label{sec:host-org}


\subsection{IBM: A Global Leader in Technology and Innovation}
\begin{figure}[H]
    \centering
    
\includegraphics[width=0.75\linewidth]{IBM.png}
    \caption{IBM}
    \label{fig:ibm-logo}
\end{figure}

\label{subsec:ibm}
IBM (International Business Machines Corporation) is an American multinational corporation located in Armonk, New York, USA, active in the field of computer hardware, software, and computer services (hosting and consulting solutions), and operating in over 170 countries.

The company was created on 1911, from the combination of the ``Computing Scale Company'' and the ``Tabulating Machine Company'' under the name of ``Computing Tabulating Recording Company'' (CTR). It changed its name to turn into International Business Machines Corporation on February 14, 1924, then recognized itself under the nickname ``Big Blue'' in reference to dark blue, a color long associated with the company and with the logo presented in the figure above.

IBM benefits from a global network of collaborators called IBMers whose mindset and values reflect the modernity of the company as well as its desire to establish a collaborative spirit. This IBM network is made up of more than 400,000 employees worldwide. Five Nobel Prizes, six Turing Awards, ten National Medals of Technology (USA), and five National Medals of Science have been awarded to IBM employees (USA).

IBM meets the needs of its customers at the international level through the implementation of its subsidiaries, branches, laboratories, and research centers in the five continents.

\begin{figure}[H]
  \centering
  \includegraphics[width=\textwidth]{ibm_imple_around_the_world.jpg}
  \caption{IBM's implementation around the world (source: IBM)}
\end{figure}

\subsection{IBM Products and Inventions}

IBM is a multinational company that manufactures and sells computer hardware, middleware, and software, as well as providing consulting services in a variety of fields, from computer systems to nanotechnology. IBM is also a big research institution, having acquired the record for the most U.S. patents generated by a company for 28 years (as of 2020). The automated teller machine (ATM), floppy disk, hard disk drive, magnetic stripe card, relational database, SQL programming language, UPC barcode, and dynamic random-access memory (DRAM) are just a few of IBM's inventions.\footnote{\url{https://en.wikipedia.org/wiki/IBM}}

IBM's business operations have consistently moved to focus on higher-value, more profitable areas. This includes the 1991 spin-off of printer manufacturer Lexmark, the 2005 and 2014 sales of Lenovo's personal computer (ThinkPad/ThinkCentre) and x86-based server operations, and the acquisition of companies like PwC Consulting (2002), SPSS (2009), The Weather Company (2016), and Red Hat (2019). IBM announced in 2015 that it would go ``fabless,'' continuing to design semiconductors but outsourcing manufacturing to GlobalFoundries, and in 2020, the company announced the spin-off of its Global Technology Services division's Managed Infrastructure Services unit, with completion expected by the end of 2021.\footnote{\url{https://newsroom.ibm.com/mergers-and-acquisitions}}

\subsection{IBM Morocco}

IBM has had an office in Morocco for 83 years, the following timeline presents the history of IBM Morocco from its first office to its last major activities in 2017.

\begin{figure}[H]
  \centering
  \includegraphics[width=\textwidth]{ibm_morocco_timeline.png}
  \caption{IBM Morocco timeline (source: IBM)}
\end{figure}

IBM Morocco operates in three major areas, see Figure~\ref{fig:ibm_morocco_services}:

\begin{itemize}
  \item \textbf{Global Business Services or GBS:} which means Business consulting. IBM Maroc markets and delivers specific development. It contains the CIC which is the cost center affiliated with the GBS activity. For example, if the GBS markets the projects, CIC delivers them at a lower cost.
  \item \textbf{Services for IT infrastructures or GTS (Global Technology Services)} which brings together a set of high value-added service offers to support companies in their transformation and meet their new global challenges. It is the specialized branch in hardware.

\begin{figure}[H]
  \centering
  \includegraphics[width=0.5\textwidth]{ibm_morocco_services.png}
  \caption{IBM Morocco Services (source: made by me)}
  \label{fig:ibm_morocco_services}
\end{figure}
\end{itemize}

\subsection{IBM Client Innovation Center}

\subsubsection{CIC Morocco}

Since January 2015, Morocco CIC has integrated MEA Centers to leverage the right skills from different centers. The acquired company MedIT in end of 2015, was integrated to IBM in 2016 and joined GIC network in 2017. Though delivering from different locations, MEA centers share the same operations, solutioning and delivery excellence functions.

The centers based in Morocco offers clients High Value application development, application maintenance and systems integration services that address increasing demand for flexible software capability, putting focus on French speaking markets in North America, Europe, and Africa.

CIC Morocco consists of a highly motivated, experienced, and dynamic workforce planning for aggressive resource growth in 2017 supported by GBS across growth platforms.

More specifically the combined entity will ensure:
\begin{itemize}
  \item Cost optimized operations through shared support resources
  \item Synchronized capacity planning execution and fulfillment optimization --- Centralized
  \item Management reporting and communication channels
  \item Common capacity build programs synchronized with LK
\end{itemize}

\subsubsection{CIC Organizational Chart}

\begin{figure}[H]
  \centering
  \includegraphics[width=\textwidth]{cic_org_chart.png}
  \caption{CIC Organizational Chart (source: IBM ``Confidential'')}
\end{figure}

Based on Casablanca, CIC Morocco is led by Mr. Mohamed Berrada and has two main services as follows:

\begin{itemize}
  \item Business Transformation Services (BTS)
  \item Hybrid Cloud Services (HCS)
\end{itemize}

The BTS is also led by Mr. Mohamed Berrada, and it is formed by several services for offering transformation consulting to the client such as Customer Transformation, Industry transformation, Data\&Technology Transformation (DTT), and Finances Supply chain Transformation (FSCT).

The DDT service line is led by Mr. Nabil RIFKI and composed of numerous data services including Data Strategy, Amazon Web Services (AWS), Data services, Microsoft, and Automation.

In the next section, we will present the Microsoft team and its organizational chart.

\subsection{Hakkoda: Snowflake Data Engineering Expertise}
\label{subsec:hakkoda}
\begin{figure}[H]
    \centering
    
\includegraphics[width=0.75\linewidth]{hakkoda.png}
    \caption{Hakkoda}
    \label{fig:hakkoda-logo}
\end{figure}

\textbf{Hakkoda} is a data engineering consultancy that built its reputation as one of the foremost Snowflake-specialized partners in North America and Europe. Founded with a focus on cloud data platform modernization, Hakkoda developed deep expertise in Snowflake architecture, \gls{dbt} pipeline development, data vault modeling, and DataOps automation. The company served a diverse portfolio of enterprise clients across the retail, manufacturing, financial services, and food and beverage sectors.

Hakkoda's service offering centered on helping organizations migrate from legacy on-premises data warehouses to Snowflake, design scalable data architectures, and implement \gls{cicd} practices for data engineering workflows. The firm distinguished itself from generalist consulting firms through its Snowflake-native engineering methodology and its emphasis on automation, repeatability, and governance.

In 2023, Hakkoda was acquired by IBM, a strategic move that allowed IBM to significantly reinforce its Snowflake practice capabilities within the CIC network. The acquisition brought Hakkoda's engineering talent, proprietary accelerators, and Snowflake Elite Partner status under the IBM umbrella, creating a combined offering that marries IBM's global delivery scale with Hakkoda's specialized cloud data engineering competencies.

The team responsible for the ARYZTA project operates under the Hakkoda brand within IBM, and is composed of senior Data Engineers, DataOps specialists, and Snowflake architects with extensive experience in enterprise data platform modernization. This internship was conducted within this team, contributing directly to the Snowflake environment alignment and CI/CD implementation workstream.

\subsection{ARYZTA: The Client Organization}
\label{subsec:aryzta}

\begin{figure}[H]
    \centering
    \placeholder{ARYZTA company logo}{Image to be inserted (file \texttt{aryzta.png}).}
    \caption{ARYZTA}
    \label{fig:aryzta}
\end{figure}


\textbf{ARYZTA AG} is a Swiss-Irish multinational corporation and one of the world's largest industrial bakery companies. Operating in more than 20 countries across Europe, North America, and Asia-Pacific, ARYZTA produces and distributes a wide range of frozen and fresh bakery products to quick service restaurants, food service operators, retail chains, and convenience stores. The company supplies products to some of the world's most recognized fast-food chains and employs tens of thousands of people globally.

ARYZTA's operations are characterized by high supply chain complexity, tight food safety and traceability requirements, and the need to coordinate logistics, production, and procurement across a geographically dispersed manufacturing network. This complexity generates large volumes of operational data that must be processed, integrated, and analyzed in near real-time to support decision-making across the business.

\subsubsection{Snowflake in the ARYZTA Data Ecosystem}

To address its data integration and analytics requirements, ARYZTA adopted \textbf{Snowflake} as its central cloud data platform. Snowflake's architecture---which decouples storage from compute, supports multi-cluster warehousing, and provides native data sharing capabilities---made it particularly well-suited to ARYZTA's needs for scalable, cost-efficient data processing across multiple business domains.

Within the ARYZTA data ecosystem, Snowflake serves as the integration layer for data sourced from ERP systems, logistics platforms, and commercial applications. Data pipelines load, transform, and expose curated datasets that feed downstream BI and reporting tools used by finance, operations, and supply chain teams. The Snowflake environment is therefore a mission-critical component of ARYZTA's operational infrastructure, and its stability directly affects the quality and availability of business-critical reporting.

The platform is structured around three separate Snowflake accounts---Development, User Acceptance Testing, and Production---each intended to correspond to a distinct stage in the software delivery lifecycle. As described in subsequent sections, the absence of governed deployment practices has compromised the integrity of this multi-environment architecture.

% ─────────────────────────────────────────────────────────────
\section{Project Presentation}
\label{sec:project-presentation}

\subsection{Context and Motivation}
\label{subsec:context-motivation}

The ARYZTA data platform was managed prior to Hakkoda's involvement by a third-party consultancy. During that period, the platform underwent extensive development, but without adherence to structured DevOps or DataOps practices. The absence of a version-controlled deployment workflow, combined with a culture of direct production modifications, resulted in a state of progressive environment drift that had accumulated over a period of several years.

When Hakkoda assumed responsibility for the platform, the team was confronted with a data infrastructure in which the three Snowflake environments---\texttt{DEV}, \texttt{UAT}, and \texttt{PROD}---had diverged to such a degree that \texttt{PROD} could no longer be considered a faithful replica of any formally documented state. Tables, stored procedures, tasks, streams, and views existed in production that had no counterpart in the lower environments, and vice versa. Schema structures were inconsistent across environments, and pipeline configurations had been individually tuned in production without those changes being propagated back to \texttt{DEV} or \texttt{UAT}.

This situation created a self-reinforcing problem: because \texttt{DEV} and \texttt{UAT} were not aligned with \texttt{PROD}, testing new features or changes in lower environments provided no reliable guarantee of behavior in production. As a result, engineers were routinely compelled to bypass the lower environments and apply changes directly to \texttt{PROD}---a practice that further exacerbated the drift and increased the probability of production incidents.

The business motivation for addressing these issues is straightforward: a stable, governed data platform is a prerequisite for reliable analytics, confident product development, and sustainable platform evolution. The implementation of environment synchronization and CI/CD automation is therefore not merely a technical improvement but a foundational investment in ARYZTA's data engineering capability.

\subsection{Existing Snowflake Environment Overview}
\label{subsec:existing-env}

The ARYZTA Snowflake landscape is organized around three independent accounts, each hosted on Snowflake's cloud infrastructure and corresponding to a standard software delivery environment:

\begin{table}[H]
  \centering
  \rowcolors{2}{lightgray}{white}
  \begin{tabularx}{\textwidth}{l l X}
    \toprule
    \rowcolor{ibmblue!15}
    \textbf{Environment} & \textbf{Identifier} & \textbf{Intended Purpose} \\
    \midrule
    Development    & \texttt{DEV}  & Feature development, unit testing, exploratory data modeling \\
    User Acceptance Testing & \texttt{UAT} & Integration testing, business validation, pre-production staging \\
    Production     & \texttt{PROD} & Live operational data, production pipelines, business-critical reporting \\
    \bottomrule
  \end{tabularx}
  \caption{ARYZTA Snowflake Environment Summary}
  \label{tab:env-summary}
\end{table}

Each environment contains a set of databases, schemas, tables, views, stored procedures, tasks, streams, and Snowflake stages that collectively constitute the data platform. In an ideal state, all three environments would be structurally synchronized, differing only in the data they contain and the compute resources allocated to their virtual warehouses.

Prior to the work undertaken in this internship, the actual state of the three environments deviated significantly from this ideal. The deployment workflow in use was informal and largely undocumented: changes were developed by engineers with direct access to each environment's SQL console, and deployments were executed through manual copy-paste of SQL scripts without version control, peer review, or systematic testing.

\begin{figure}
    \centering
    \includegraphics[width=1\linewidth]{ci_cd.png}
    \caption{Before CI/CD}
    \label{fig:placeholder}
\end{figure}

\subsection{Risks and Challenges}
\label{subsec:risks-challenges}

The combination of environment drift and ungoverned deployment practices exposes ARYZTA's data platform to a range of technical and operational risks. These may be categorized along two dimensions: deployment risk and maintenance risk.

\paragraph{Deployment Risk.} The absence of a validated lower-environment proxy for production means that the consequences of any given deployment are difficult to anticipate. A change deployed directly to \texttt{PROD} without prior testing may introduce schema conflicts, break downstream transformations, or corrupt dependent pipelines. Given the business-critical nature of the ARYZTA data platform, such incidents can directly affect operational reporting and supply chain visibility.

\paragraph{Maintenance and Auditability Risk.} Without version control, there is no definitive record of what changes were applied to each environment, by whom, and when. This makes incident debugging significantly more difficult, as engineers cannot reliably reconstruct the sequence of changes that preceded a failure. It also complicates the onboarding of new team members, who must infer the current state of the platform through direct environment inspection rather than consulting a versioned codebase.

\paragraph{Governance Risk.} The lack of \gls{rbac}-enforced deployment workflows means that any engineer with Snowflake account access can modify production objects directly. This represents a significant governance exposure, particularly in the context of data compliance requirements.

\paragraph{Scalability Risk.} As the platform evolves and the number of Snowflake objects grows, the complexity of maintaining consistency between environments through manual means increases non-linearly. Without automation, the technical debt associated with environment misalignment will continue to accumulate.

\subsection{Specific Problem Addressed}
\label{subsec:specific-problem}

Based on the preliminary assessment conducted at the outset of the internship, the specific technical problem addressed in this project can be formally defined as follows:

\begin{tcolorbox}[notebox, title={Problem Statement}]
ARYZTA's Snowflake data platform currently operates across three environments---DEV, UAT, and PROD---that have diverged significantly in schema structure, object definitions, and pipeline configuration due to years of ungoverned, manual, and undocumented deployments. The absence of a version-controlled deployment process and an automated CI/CD pipeline prevents the engineering team from safely and reproducibly promoting changes through the environment lifecycle, forcing direct production deployments and continuously increasing the risk of data platform failures.
\end{tcolorbox}

This problem encompasses the following specific technical dimensions:

\begin{enumerate}
  \item \textbf{Environment Synchronization}: DEV and UAT do not accurately reflect the PROD state, making lower-environment testing unreliable.
  \item \textbf{Missing Snowflake Objects}: Numerous tables, stored procedures, tasks, and streams exist in PROD that are absent from DEV and/or UAT.
  \item \textbf{Inconsistent Schema Definitions}: Object definitions (column types, constraints, default values) differ across environments, introducing transformation failures when deployments are attempted.
  \item \textbf{Absent Deployment Automation}: There is no automated pipeline to validate, test, and promote SQL-based changes across environments in a controlled manner.
  \item \textbf{Lack of Version Control}: No central repository tracks the authoritative state of the Snowflake codebase, making the current platform state non-reproducible.
\end{enumerate}

\subsection{Project Objectives}
\label{subsec:objectives}

In response to the identified problem, the following project objectives were formally agreed upon with the Hakkoda engineering team and the ARYZTA client stakeholders:

\begin{table}[H]
  \centering
  \rowcolors{2}{lightgray}{white}
  \begin{tabularx}{\textwidth}{c l X}
    \toprule
    \rowcolor{ibmblue!15}
    \textbf{\#} & \textbf{Objective} & \textbf{Description} \\
    \midrule
    O1 & Environment Audit       & Systematically document all discrepancies between DEV, UAT, and PROD across all Snowflake object types. \\
    O2 & Root Cause Analysis     & Identify and characterize the technical and organizational factors that led to the observed environment drift. \\
    O3 & Synchronization Strategy & Design a phased approach to bring DEV and UAT into alignment with PROD. \\
    O4 & CI/CD Pipeline          & Implement an automated deployment pipeline integrated with GitHub that enforces version-controlled promotions. \\
    O5 & GitHub Integration      & Establish a Git-based development workflow with branch protection rules and pull request review gates. \\
    O6 & Governance Framework    & Define and document deployment policies, role assignments, and environment access controls. \\
    O7 & Documentation           & Produce comprehensive technical documentation covering the architecture, the pipeline, and the onboarding process. \\
    \bottomrule
  \end{tabularx}
  \caption{Project Objectives Summary}
  \label{tab:objectives}
\end{table}

% ─────────────────────────────────────────────────────────────
\section{Project Management}
\label{sec:project-mgmt}

\subsection{Agile Methodology: SCRUM}
\label{subsec:scrum}

The project was managed in accordance with the \textbf{Agile SCRUM} framework, a widely adopted iterative methodology in software and data engineering contexts. SCRUM's emphasis on short delivery cycles, continuous feedback, and adaptive planning made it particularly well-suited to the nature of this project, where the full scope of discrepancies between environments could only be determined progressively as the investigation advanced.

The SCRUM workflow was structured around the following recurring ceremonies and artefacts:

\begin{itemize}
  \item \textbf{Sprint Planning}: At the beginning of each two-week sprint, the team reviewed the backlog and selected a set of user stories and technical tasks to be completed within the sprint. Effort estimation was conducted using story points.
  \item \textbf{Daily Stand-ups}: Short daily synchronization meetings (15 minutes) allowed team members to communicate progress, identify blockers, and coordinate cross-functional dependencies.
  \item \textbf{Sprint Reviews}: At the end of each sprint, completed deliverables were demonstrated to stakeholders, providing an opportunity for feedback and course correction.
  \item \textbf{Sprint Retrospectives}: The team reflected on the process itself---identifying what worked well, what could be improved, and what actions would be taken in the next sprint.
  \item \textbf{Backlog Refinement}: The Product Backlog was continuously groomed to ensure that upcoming stories were well-defined, estimated, and prioritized.
\end{itemize}

\subsubsection{Tooling and Collaboration Infrastructure}

The project made use of the following tools to support SCRUM execution:

\begin{table}[H]
  \centering
  \rowcolors{2}{lightgray}{white}
  \begin{tabularx}{\textwidth}{l l X}
    \toprule
    \rowcolor{ibmblue!15}
    \textbf{Tool} & \textbf{Category} & \textbf{Usage} \\
    \midrule
    ServiceNow & ITSM / Ticketing       & Sprint backlog, task tracking, incident management, change request documentation \\
    GitHub     & Version Control        & Source code repository, branch management, pull request workflow, CI/CD pipeline triggers \\
    GitHub Actions & CI/CD Automation   & Automated deployment pipeline execution, environment promotion workflows \\
    Slack      & Communication          & Team communication, sprint updates, stakeholder notifications \\
    Confluence & Documentation          & Technical documentation, architecture records, runbooks \\
    \bottomrule
  \end{tabularx}
  \caption{Project Tooling and Collaboration Infrastructure}
  \label{tab:tools}
\end{table}

The use of \textbf{ServiceNow} as the primary ticketing platform is noteworthy, as it provided traceability between individual technical tasks and the broader change management process required by ARYZTA's internal IT governance policies. Every change deployed to the PROD environment was required to be associated with a ServiceNow change request, ensuring auditability.

\subsection{Project Timeline}
\label{subsec:gantt}

The internship was structured into six sequential phases, each corresponding to a distinct set of deliverables and outcomes. The phases are described below and visualized in the Gantt chart that follows.

\begin{enumerate}
  \item \textbf{Onboarding and Environment Discovery} (Weeks 1--2): Familiarization with the Snowflake platform, IBM/Hakkoda tooling, and the ARYZTA project context. Initial access provisioning and documentation review.
  \item \textbf{Environment Audit and Discrepancy Analysis} (Weeks 3--5): Systematic comparison of DEV, UAT, and PROD environments across all object types. Documentation of all identified discrepancies in a structured audit report.
  \item \textbf{Root Cause Analysis and Strategy Design} (Weeks 6--8): Investigation of the origins of environment drift. Design of the synchronization and CI/CD strategy, including architecture decisions and tooling selection.
  \item \textbf{CI/CD Pipeline Implementation} (Weeks 9--14): Development of the automated deployment pipeline using GitHub Actions, including scripting of validation, deployment, and rollback procedures.
  \item \textbf{Environment Synchronization Execution} (Weeks 15--17): Phased application of the synchronization strategy to bring DEV and UAT into alignment with PROD.
  \item \textbf{Testing, Documentation, and Handover} (Weeks 18--20): End-to-end pipeline testing, production of technical documentation, and knowledge transfer to the ARYZTA and Hakkoda teams.
\end{enumerate}

\newpage
\begin{landscape}
\begin{figure}[H]
  \centering
  \begin{ganttchart}[
      hgrid,
      vgrid={*{6}{draw=none}, dotted},
      x unit=0.55cm,
      y unit title=0.7cm,
      y unit chart=0.65cm,
      title height=1,
      title label font=\small\bfseries,
      bar/.style={fill=ibmblue!70, rounded corners=1pt},
      bar label font=\small,
      group label font=\small\bfseries,
      milestone label font=\small,
      link/.style={->, thick, accentblue},
      bar height=0.5,
    ]{1}{20}
    \gantttitle{Internship Timeline (Weeks)}{20} \\
    \gantttitlelist{1,...,20}{1} \\

    \ganttgroup{Phase 1: Onboarding}{1}{2} \\
    \ganttbar{Environment discovery}{1}{2} \\
    \ganttbar{Tool access setup}{1}{2} \\

    \ganttgroup{Phase 2: Audit}{3}{5} \\
    \ganttbar{DEV/UAT/PROD comparison}{3}{5} \\
    \ganttbar{Discrepancy documentation}{4}{5} \\

    \ganttgroup{Phase 3: Analysis \& Design}{6}{8} \\
    \ganttbar{Root cause analysis}{6}{7} \\
    \ganttbar{Architecture design}{7}{8} \\

    \ganttgroup{Phase 4: CI/CD Implementation}{9}{14} \\
    \ganttbar{GitHub repo setup}{9}{9} \\
    \ganttbar{Pipeline scripting}{9}{13} \\
    \ganttbar{Pipeline testing}{12}{14} \\

    \ganttgroup{Phase 5: Synchronization}{15}{17} \\
    \ganttbar{DEV sync execution}{15}{16} \\
    \ganttbar{UAT sync execution}{16}{17} \\

    \ganttgroup{Phase 6: Docs \& Handover}{18}{20} \\
    \ganttbar{Technical documentation}{18}{19} \\
    \ganttbar{Knowledge transfer}{19}{20} \\
    \ganttmilestone{Final delivery}{20}

  \end{ganttchart}
  \caption{Project Gantt Chart --- Internship Timeline}
  \label{fig:gantt}
\end{figure}
\end{landscape}

\section*{Chapter Conclusion}

This chapter has established the organizational and technical foundation upon which the remainder of this report is built. The host organizations---IBM CIC Morocco and Hakkoda---bring complementary capabilities to the ARYZTA engagement: IBM's global delivery infrastructure and enterprise credibility, and Hakkoda's Snowflake-native engineering expertise. The client, ARYZTA, operates a mission-critical Snowflake data platform that has suffered from years of ungoverned deployment practices, resulting in a state of environment drift that constrains the engineering team's ability to develop and deliver platform enhancements safely.

The project objectives are clearly defined: to audit and resolve the environment discrepancies, to implement a CI/CD pipeline that enforces version-controlled deployments, and to establish a governance framework that will prevent a recurrence of the current situation. The Agile SCRUM methodology, supported by ServiceNow and GitHub, provides the structured delivery framework within which these objectives are pursued.

The following chapter presents a rigorous analysis of the current state of the platform, including a detailed characterization of the discrepancies observed, their operational impact, and the factors that gave rise to them.

\newpage
